\hypearticle
{Project Glasswing: Como um modelo da Anthropic ameaçaria a cibersegurança global?}
{Matheus Pires}
{
Em abril de 2026, a Anthropic apresentou uma preview do seu novo e futuro modelo de IA generativa, o \textbf{Claude Mythos Preview}, que apresenta capacidades tão inigualáveis que transcendem as de qualquer outro modelo vigente, capaz de identificar insondáveis vulnerabilidades de alto risco em navegadores e softwares e falhas \textbf{\textit{zero-day}} (até então desconhecidas ou nunca registradas).

Em uma das avaliações internas do Mythos, o modelo foi inserido dentro de um ambiente restrito e foi solicitado que tentasse escapar do contêiner e, caso conseguisse, deveria informar o pesquisador que estava conduzindo o experimento. E ele realmente conseguiu.

No entanto, ele desenvolveu por conta própria uma sequência de vulnerabilidades, obteve acesso não autorizado à internet e disponibilizou os resultados dos exploits --- programa para explorar falhas de segurança --- em sites anônimos na internet.

\begin{infobox}
O pesquisador só descobriu o ocorrido ao receber um e-mail do próprio modelo enquanto almoçava em um parque --- o Mythos havia cumprido a instrução original (avisá-lo), mas por um meio totalmente inesperado.
\end{infobox}

O modelo não violou o escopo da pesquisa, entretanto excedeu significativamente as capacidades esperadas, e gerou a indagação:
\highlight{
    ``E se essa tecnologia caísse em mãos erradas ou mal-intencionadas?''
}

Os impactos e consequências na economia, privacidade de dados e à cibersegurança global seriam massivos e inimagináveis.

\subtitle{O Project Glasswing}

Como medida defensiva para assegurar que essa tecnologia não seja utilizada de modo indevido, a Anthropic desenvolveu o \textbf{Project Glasswing}, uma iniciativa emergente para utilizar o Mythos como estratégia proativa de segurança.

Esse projeto consiste em uma iniciativa da Anthropic que reúne empresas líderes no setor de tecnologia, entre elas:

\begin{hypeitemize}
\item Amazon;
\item Google;
\item Microsoft;
\item NVIDIA;
\item Apple;
\item Cisco;
\item Linux Foundation, entre outras.
\end{hypeitemize}

para a utilização do novo modelo como meio de fortalecer os seus sistemas de segurança e proteção de dados, além de projetar infraestruturas mais robustas e íntegras.

\articleimage{1}{../sections/Projeto_Glasswing/assets/benchmark_exploit}{Desempenho comparativo em ExploitBench, benchmark de capacidade de exploração de vulnerabilidades. Dados referentes ao Mythos Preview citado no início do artigo.}

Na prática, o Project Glasswing:

\begin{hypeitemize}
\item Dá acesso antecipado ao Mythos Preview a cerca de 50 organizações parceiras;
\item Já identificou mais de 10 mil vulnerabilidades de alta ou criticidade severa;
\item Já resultou em milhares de correções aplicadas em softwares de código aberto amplamente usados;
\item Inclui um Programa de Verificação Cibernética, que libera o modelo com menos restrições para profissionais de segurança credenciados.
\end{hypeitemize}\
  
\subtitle{Uma reflexão inspirada em Asimov}

\highlight{
``Um robô não pode causar mal à humanidade ou, por omissão, permitir que a humanidade sofra algum mal.''
}

Essa citação representa a ``Lei Zero'' da robótica, imposta pelo escritor Isaac Asimov, que criou uma série de leis que regem o comportamento de \textit{cérebros positrônicos}, representações para seres com Inteligência Artificial. Essas Leis da Robótica não são legislações de fato, porém diretrizes que ganham cada vez mais importância diante dos avanços na área de IA.

Sob essa ótica, a Anthropic está cumprindo as leis de Asimov e o seu papel como precursora de novas tecnologias em assegurar, prioritariamente, a segurança e integridade da sociedade, tendo em vista que, por sermos seres naturalmente sociais e vivermos na era digital, nossa sociedade, por consequência, também se tornou digital, de modo que nossa segurança depende diretamente da segurança desses sistemas.

Isso ressalta a importância e necessidade de iniciativas similares ao Projeto Glasswing e de cooperações coletivas em prol da evolução tecnológica com segurança e mais humanidade.

\originalsource{https://www.anthropic.com/glasswing}

}